\originalsection{作业 4: 约束与阻力}
\begin{exercise}
原题 1. 倾角 $\theta$ 的斜面上块 $m_1$, 绳连该块后越过固定轮, 绕承托 $m_2$ 的无质量动轮, 最后固定; 动轮两侧绳竖直. 块与斜面静摩擦 $\mu_s$. (a) 求使块刚向上滑的最小 $m_2$. (b) 求使块刚向下滑的最大 $m_2$. (c) 若块正在上滑, 动摩擦 $\mu_k$, 求块加速度.
\end{exercise}
\begin{solution}
静止动轮给 $T=m_2g/2$. (a) 上滑边界 $T=m_1g(\sin\theta+\mu_s\cos\theta)$, 所以 $m_{2,\min}=2m_1(\sin\theta+\mu_s\cos\theta)$. (b) 下滑边界 $m_{2,\max}=2m_1(\sin\theta-\mu_s\cos\theta)$, 若该量为负, 任何非负 $m_2$ 都不能使静止块向下起滑. (c) 取块向上、动轮向下为正, 绳长给 $a_1=2a_2$. 联立 $T-m_1g(\sin\theta+\mu_k\cos\theta)=m_1a_1$ 与 $m_2g-2T=m_2a_1/2$, 得
\begin{equation}
a_1=\frac{2g[m_2-2m_1(\sin\theta+\mu_k\cos\theta)]}{4m_1+m_2}.
\end{equation}
负号表示上滑过程中减速, 原题若问大小应取绝对值, 不能改动摩擦方向.
\end{solution}
\begin{exercise}
原题 2. 质量 $m$ 粒子起初沿 $+x$ 以 $v_0$ 运动, 受反向力 $\vct F=-be^{cv}\uv x$, $b>0$ 单位 N, $c>0$ 单位 s/m. 求速率 $v(t)$.
\end{exercise}
\begin{solution}
$me^{-cv}\dd v=-b\dd t$, 积分得 $e^{-cv}=e^{-cv_0}+bc t/m$, 因而 $v=-(1/c)\ln(e^{-cv_0}+bc t/m)$. 该式只到首次停止 $t_*=m(1-e^{-cv_0})/(bc)$; 以后“反向力”的方向未由 $-\uv x$ 原式继续正确定义, 不能把负数作为速率.
\end{solution}
\begin{exercise}
原题 3. 均匀绳质量 $m_1$、长 $l$, 一端连转轴, 另一端连点质量 $m_2$. 系统以角速度 $\omega$ 近似水平旋转, 忽略重力及轴半径. 求距轴 $r$ 处张力.
\end{exercise}
\begin{solution}
截面须提供外侧点质量和绳段的向心力, $T(r)=m_2\omega^2l+\int_r^l(m_1/l)\omega^2s\dd s=m_2\omega^2l+m_1\omega^2(l^2-r^2)/(2l)$. 末端边界 $T(l)=m_2\omega^2l$.
\end{solution}
\begin{exercise}
原题 4. 轻绳跨固定圆杆半周, 两端悬 $m_1>m_2$. 静摩擦 $\mu_s$, 动摩擦 $\mu_k<\mu_s$. 重端切点 A, 轻端切点 B, 接触角 $\pi$, 极限关系 $T_B=T_Ae^{-\mu\pi}$. (a) 比较 $T_A$ 与重块所受张力 $T_1$. (b) 求开始滑动的重质量阈值. (c) 已滑动时求加速度.
\end{exercise}
\begin{solution}
(a) 无质量自由段张力相同, $T_A=T_1$. (b) 静止阈值 $m_1=m_2e^{\mu_s\pi}$. (c) 令 $q=e^{-\mu_k\pi}$, 重块向下 $m_1g-T_A=m_1a$, 轻块向上 $qT_A-m_2g=m_2a$. 消去 $T_A$ 得 $a=g(qm_1-m_2)/(qm_1+m_2)$.
\end{solution}
\begin{exercise}
原题 5. 质量 $m$ 的球以水平初速 $u\uv x$ 射入液体, 初始竖直速度为零. 阻力 $-C\vct v$, 向下为 $+y$. (a)(b) 写两方向加速度. (c)(d) 求 $v_x(t),v_y(t)$. (e)(f) 求两个终端分量.
\end{exercise}
\begin{solution}
(a) $\dot v_x=-Cv_x/m$. (b) $\dot v_y=g-Cv_y/m$. 令 $\tau=m/C$, 积分初值给 (c) $v_x=ue^{-t/\tau}$, (d) $v_y=g\tau(1-e^{-t/\tau})$. (e) $v_{xT}=0$, (f) $v_{yT}=mg/C$. 向下正轴决定重力符号.
\end{solution}
\begin{exercise}
原题 6. 上绳经过固定天花板轮 A, 一端接 $m_1$, 另一端接动轮 B 的轴. 下绳一端固定在地板, 向上绕 B 后另一端接 $m_2$. 绳轮均无质量不可伸长, 求 $m_2$ 加速度, 向下为正.
\end{exercise}
\begin{solution}
上绳给 $a_B=-a_1$, 下绳接地使 $a_2=2a_B=-2a_1$. 设下绳张力 $T$, B 平衡给上绳张力 $2T$. 块方程 $m_1g-2T=m_1a_1$, $m_2g-T=m_2a_2$. 代约束得 $a_2=2g(2m_2-m_1)/(m_1+4m_2)$. 固定在地板这一边界不可改成天花板.
\end{solution}
